\section*{الفصل \ref{ch:g8:speed-of-light} --- سرعة الضوء؛ المسافات في الكون}

\begin{solution}{exo:g8:speed-of-light:1}
$c = \qty{300000}{km/s} = 3 \times 10^{8}\ \unit{m/s}$.
وفانوسا التلّين لم يقيسا إلا ردود الفعل البشرية: فقد عبر
\omterm{def:g1:light-and-shadows:source}{الضوء} بين التلّين في أجزاء من مليون من الثانية، أي أدنى بكثير
من توقيت أيّ يد.
\end{solution}

\begin{solution}{exo:g8:speed-of-light:2}
\omterm{def:g2:moon-in-the-sky:moon}{قمر} المشتري ينخسف بجدول منضبط --- ساعة في السماء. و\emph{خبر}
كل \omterm{def:g7:shadows-eclipses:lunar}{خسوف} كان يصل متأخّرًا كلما وقفت الأرض بعيدًا عن المشتري،
ومتقدّمًا كلما قربت. والتأخّر قاس زمن سفر \omterm{def:g1:light-and-shadows:source}{الضوء} عبر عرض \omterm{def:g6:sun-earth-moon:axisorbit}{مدار}
الأرض --- ومن ثمّ \omterm{def:g7:motion-average-speed:formula}{سرعة} \omterm{def:g1:light-and-shadows:source}{الضوء}.
\end{solution}

\begin{solution}{exo:g8:speed-of-light:3}
\omterm{def:g2:moon-in-the-sky:moon}{القمر}: $384000 \div 300000 \approx \qty{1.3}{s}$. والشمس:
$150000000 \div 300000 = \qty{500}{s}$ --- أي ثماني دقائق
وعشرون ثانية. \omterm{def:g1:light-and-shadows:source}{فضوء} الشمس يُري الشمس دائمًا كما كانت قبل ثماني
دقائق: خبر قديم بحكم البناء.
\end{solution}

\begin{solution}{exo:g8:speed-of-light:4}
مسافة: هي ما يقطعه \omterm{def:g1:light-and-shadows:source}{الضوء} في سنة واحدة --- نحو
$9.5 \times 10^{12}$ كيلومتر.
\end{solution}

\begin{solution}{exo:g8:speed-of-light:5}
أربع سنوات. والجواب يختلف من قارئ إلى آخر --- فقبل أربع سنوات
كنت على بعد صفوف قليلة إلى الوراء في هذا الكتاب بعينه.
\end{solution}

\begin{solution}{exo:g8:speed-of-light:6}
عند نحو إحدى عشرة دقيقة في الاتجاه الواحد، يجيب أيّ تصحيح
توجيه عن صورة عمرها اثنتان وعشرون دقيقة على الأقلّ --- فتكون
المركبة في الصدع قبل أن تصل رجفة عصا التوجيه. ومن هنا الخطط،
لا عصيّ التوجيه.
\end{solution}

\begin{solution}{exo:g8:speed-of-light:7}
\omterm{def:g2:moon-in-the-sky:moon}{القمر}: نحو \qty{1.3}{s}. والشمس: ثماني دقائق. وأقرب \omterm{def:g4:solar-system:star}{نجم}: أربع
سنوات. والمجرّة المجاورة: مليونان ونصف من السنين.
\end{solution}

\begin{solution}{exo:g8:speed-of-light:8}
\omterm{def:g1:light-and-shadows:source}{الضوء}. فقد ثبتت \omterm{def:g7:motion-average-speed:formula}{السرعة} $c$ أرسخ من أيّ قضيب معدني أو مسطرة
وطنية --- فصارت المسطرة تُشتقّ من \omterm{def:g1:light-and-shadows:source}{الضوء} الآن، لا \omterm{def:g1:light-and-shadows:source}{الضوء} يُقاس
بالمسطرة.
\end{solution}

\begin{solution}{exo:g8:speed-of-light:9}
الذهاب والإياب: $300000 \times 2.6 = 780000$ كيلومتر؛ \omterm{def:g2:moon-in-the-sky:moon}{والقمر}
يقف على النصف --- \qty{390000}{km}. وقاعدة السونار تنسحب
كاملة: فكل صدى يدفع ثمن الذهاب والإياب، فنصِّف قبل أن تصدّق.
\end{solution}

\begin{solution}{exo:g8:speed-of-light:10}
\omterm{def:g3:first-electric-circuit:bulb}{المصباح}: $3 \div (3 \times 10^{8}) = 10^{-8}$ ثانية --- أي
جزء من مئة من المليون من الثانية: فحتى الغرفة أرشيف، لكنه
أضحل من أن يُلحظ. ثم الشمس (ثماني دقائق)، \omterm{def:g4:solar-system:star}{فالنجم} (قرن)،
فالمجرّة (مليونان ونصف من السنين) --- فكلما نظرت أعمق كانت
الصفحة أقدم.
\end{solution}

\begin{solution}{exo:g8:speed-of-light:11}
``يوم الثلاثاء الماضي \emph{بلغنا} \omterm{def:g1:light-and-shadows:source}{ضوء} انفجار \omterm{def:g4:solar-system:star}{نجم} --- أمّا
الانفجار نفسه فوقع قبل نحو $8000$ سنة، وخبره يسافر منذ
ذلك الحين.''
\end{solution}

\begin{solution}{exo:g8:speed-of-light:12}
$3 \times 10^{8}\ \text{km} \div 1000\ \text{s} = 3 \times
10^{5}$ \unit{km/s} --- أي القيمة الحديثة تمامًا (وقد وقع
التقدير التاريخي، بأرقام \omterm{def:g6:sun-earth-moon:axisorbit}{مدار} أخشن، في حدود ربع منها: وهو
جواب أول مذهل بمقاييس القرن السابع عشر).
\end{solution}

\begin{solution}{pb:g8:speed-of-light:1}
\textbf{1.} بعد نحو \qty{1.3}{s} من كلامك؛ وأقرب إقرار يعود
بعد رحلة الذهاب والإياب، أي نحو \qty{2.6}{s}.
\textbf{2.} $36000 \div 300000 = \qty{0.12}{s}$ في الاتجاه
الواحد؛ والسؤال والجواب يركبان صعودًا وهبوطًا مرتين --- أي
قرابة نصف ثانية من التردّد، عند عتبة الأذن تمامًا.
\textbf{3.} ``اندلع التوهّج قبل ثماني دقائق وعشرين ثانية من
إضاءة شاشاتنا --- ونحن نشاهده كما \emph{كان}.''
\textbf{4.} وصلة \omterm{def:g2:moon-in-the-sky:moon}{القمر}: \qty{1.3}{s} في مقابل \qty{0.12}{s}
--- أي نحو عشرة أضعاف الهامش.
\textbf{5.} $1.8 \times 10^{8} \div (3 \times 10^{5}) = 600$
ثانية: أي عشر دقائق في الاتجاه الواحد.
\textbf{6.} الذهاب والإياب عشرون دقيقة: فصورة الإنقاذ عمرها
عشر دقائق، وأمر الإنقاذ متأخّر عشر دقائق --- فيفوز الصدع
بتسع عشرة دقيقة. فوجب أن تحمل المركبة ردود فعلها هي: توقّفًا
ذاتيًّا عند الأخطار على متنها.
\textbf{7.} يجب أن يصل آخر جزء عند 21:00؛ فيغادر الأرض عند
$21{:}00 - 10$ دقائق $= 20{:}50$؛ وقد بدأ الإرسال قبل ذلك
بثلاث دقائق: فالبدء عند 20:47.
\textbf{8.} يتضاعف التأخير إلى نحو إحدى وعشرين دقيقة في
الاتجاه الواحد: فليس للكواكب ``مسافة'' واحدة --- فكلا
العالمين يدور، وكل محادثة تُجدوَل في مواجهة هدف متحرّك.
\textbf{9.} $300000 \times 86400 \approx 2.6 \times 10^{10}$
كيلومتر --- أي ستّة وعشرون ألف مليون كيلومتر؛ والسؤال والجواب
يستغرقان يومين كاملين.
\textbf{10.} يبلغ \omterm{def:g1:light-and-shadows:source}{الضوء} \omterm{def:g4:solar-system:star}{النجم} في نحو أربع سنوات؛ وأقرب ردّ
يحطّ بعد نحو ثماني سنوات من إرسال الليلة.
\textbf{11.} ``غادر هذا \omterm{def:g1:light-and-shadows:source}{الضوء} مجرّته قبل مليونين ونصف من
السنين: فاللقطة تبثّ أرض أوائل الأدوات الحجرية --- أعمق صفحة
في الأرشيف تبلغها العين المجرّدة.''
\textbf{12.} مثلًا: ``لا شيء لحظي --- فكل إشارة، ومنها
إشاراتنا، تسافر \omterm{def:g7:motion-average-speed:formula}{بسرعة} $c$ في أحسن الأحوال. وعصيّ التوجيه تخسر
أمام التأخير فيما وراء \omterm{def:g2:moon-in-the-sky:moon}{القمر}: فعلى الآلات البعيدة أن تفكّر
لنفسها. وكل مقراب قارئ أرشيف --- فكلما بعد الجسم قدم الخبر،
ولا طبعة أنضر يمكن طبعها.''
\end{solution}
